风险管理的目标不是罗列所有坏消息，而是提前定义什么指标会改变盈利分母、估值倍数与动作。当前最高优先级是利润率不达预期、H1盈利前置、交付/客户履约和现金转化；集团项目误归属与军品透明度则通过零计价控制。

\section{E16｜风险矩阵：概率与影响分开看}

概率档是House主观研究区间，不是基于足够历史样本的统计频率：低为\(p<25\%\)，中低为\(25\%\leq p<40\%\)，中为\(40\%\leq p<60\%\)，中高为\(60\%\leq p<75\%\)，高为\(p\geq75\%\)。影响档按对基础每股价值或投资主张的潜在损害划分：中约低于10\%，中高约10\%--20\%，高为高于20\%或足以触发硬失效。无法从冻结证据估计金额的风险明确标为\texttt{not estimable}；集团、军品和政策期权维持\texttt{zero-credit}。

\begin{exhibitbox}[风险热度与估值传导]
\scriptsize
\begin{tabularx}{\exhibitboxwidth}{L{2.3cm}C{1.2cm}C{1.2cm}L{3.5cm}L{4.0cm}X}
\toprule
\textbf{风险} & \textbf{概率} & \textbf{影响} & \textbf{硬阈值} & \textbf{财务/估值敏感性} & \textbf{动作} \\
\midrule
高价订单毛利低于预期 & 中 & 高 & 核心GM低于10.5\% & 在确认0.96下，Growth GM 15.8\%/17.8\%对应2026归母159.48/186.04亿元；14x价值较Base约-2.47/+2.47元 & 下调盈利与倍数 \\
H1盈利前置 & 中高 & 高 & H1低于92/90亿元，或Q3利润率断裂 & 低端下H2仍需80.76亿元；若持续性破坏则转Bear，2027价值11.21元 & 禁止年化，重估持续期 \\
交付/客户履约 & 中低 & 高 & 正式下调168艘/1,650万DWT，或有效交付降至约1,550万DWT & Bear 2026收入1,652.49亿元，较Base低122.51亿元；价值与倍数同步下修 & 延后收入、提高折价 \\
新船价格/订单回落 & 中 & 中高 & 价格与覆盖年限持续下行；单点门槛\texttt{not estimable} & 影响2028后利润/终值，幅度\texttt{not estimable}；未签订单不冲减已签合同 & 降低持续期 \\
现金转化恶化 & 中 & 高 & H1 OCF为负，或全年OCF/NP低于0.8x且无回收路径 & 直接去除现金质量倍数信用；每股幅度\texttt{not estimable}至完整现金桥更新 & 下调倍数 \\
汇率与套保 & 中 & 中 & 汇兑/衍生损失重大；敞口阈值\texttt{not estimable} & 合同币种和套保不足，单独敏感性\texttt{not estimable} & 下调净利率 \\
重组整合 & 中低 & 高 & 审计、关联交易、排产或少数股东重大异常 & 未量化协同溢价本来即\texttt{zero-credit}；异常则下调归属和倍数 & 取消协同信用 \\
地缘政治/制裁 & 中 & 高 & 客户、设备、融资或保险实质受限 & 项目级金额\texttt{not estimable}；政策期权\texttt{zero-credit} & 下调转化率 \\
集团项目误归属 & 中高 & 高 & 只有集团新闻、无600150公告 & 未分配集团订单、沪东中华均\texttt{zero-credit} & 维持0 \\
军品信息不透明 & 高 & 中 & 订单/收入/毛利继续未披露 & 军品收入、EPS与SOTP均\texttt{not estimable / zero-credit} & 禁止军品SOTP \\
\bottomrule
\end{tabularx}
\sourcenote{风险框架；公司阈值连接CF-03--CF-05、CF-17与冻结模型。}
\end{exhibitbox}

风险矩阵的“所以呢”是把重点放在可观察的领先指标：核心毛利、交付、OCF与营运资金。宏大叙事风险若没有公司层证据，先用零计价控制；经营风险一旦触发，则立即通过模型下修，而不是等待故事修复。

\section{催化剂日历：什么事件真正改变价值}

正式半年报是最近的关键节点，因为它首次把H1利润预告与收入、毛利、现金和资产负债表联系起来。之后依次是Q3利润持续性、全年交付与OCF、新签订单的价格/结构，以及任何正式资产或重大合同公告。

\begin{exhibitbox}[催化剂与证据要求]
\small
\begin{tabularx}{\exhibitboxwidth}{L{2.8cm}L{4.2cm}L{4.0cm}X}
\toprule
\textbf{节点} & \textbf{正向催化} & \textbf{负向验证} & \textbf{动作} \\
\midrule
2026H1正式报告 & 归母/扣非高于101/99，核心GM至少11.72\%，OCF正 & 低于92/90、GM低于10.5\%、现金背离 & 三重验证后重跑估值 \\
2026Q3 & H2利润率保持、交付正常 & H1明显前置且无节奏解释 & 调整持续期与倍数 \\
2026全年经营 & 168艘/1,650万DWT计划接近完成 & 重大延期、撤单或客户融资问题 & 更新收入确认与Bear概率 \\
季度/年度订单 & 高价值结构与价格可审计 & 价格/结构恶化、覆盖收缩 & 更新2028后利润中枢 \\
成本/汇率 & 设备、钢材与汇率被签价和效率吸收 & 设备涨价、返工与汇兑侵蚀 & 更新Growth GM \\
资产/重大合同公告 & 600150正式披露范围、价格、盈利、股本 & 只有集团新闻或媒体传闻 & 正式公告前0；公告后完整重跑 \\
\bottomrule
\end{tabularx}
\sourcenote{House view、variant perception与风险框架。}
\end{exhibitbox}

催化剂并不必然等于股价正向。H1“符合预告”但仅靠一次性项目、现金恶化或毛利不升，可能只是兑现市场基准；只有结果改变2027--2028盈利持续期，才改变内在价值。

\section{升级、维持、降级：把标签翻译成行动}

当前“中性偏多”不是模糊措辞，而是两个条件的组合：经营方向偏正面，估值安全边际不足。已有仓位可以保留事件观察，新增资金等待验证；若失效条件发生，应先降低研究敞口，不能以“长期造船强国”维持旧目标。

\begin{exhibitbox}[投委会动作规则]
\small
\begin{tabularx}{\exhibitboxwidth}{L{2.4cm}L{4.5cm}L{4.2cm}X}
\toprule
\textbf{动作} & \textbf{必要条件} & \textbf{估值条件} & \textbf{执行含义} \\
\midrule
升级 & 利润、核心GM、OCF三重验证；交付接近计划 & 独立Base相对现价至少15\%上行，压力下行可控 & 进入新增仓位评估 \\
维持 & 业绩在预告区间但质量或持续性仍不清晰 & 公允价值与现价大体重合 & 持有等待；不因单一指标追价 \\
降级 & H1低于92/90、GM低于10.5\%、重大延期/撤单或现金恶化 & 盈利分母或倍数下修 & 减少观察敞口并重跑模型 \\
退出观察 & 主体边界、治理或重大风险改变可投资前提 & Bear仍不能覆盖风险 & 终止旧主张，等待新证据包 \\
\bottomrule
\end{tabularx}
\sourcenote{House view与冻结估值规则；不构成证券交易指令。}
\end{exhibitbox}

“至少15\%上行”是升级研究门槛，不是收益承诺。即便正式结果强，如果股价已同步上涨导致新公允价值相对现价仍无安全边际，动作也可能维持；相反，价格下跌只有在失效条件未触发时才改善赔率。

\section{最强反方与双向证伪}

最强反方认为等待现金确认过于保守：高在手、交付计划和低远期PE已使盈利上修确定性较高，等证据齐全再买会损失回报。证明我们过于谨慎的组合证据是：H1达到或超过110/108亿元，核心毛利显著改善，OCF同步增强，2026交付按计划，并披露更清晰的2027--2028高价订单排程。

证明市场多头错误的组合证据则是：H1低于92/90亿元或依赖非经常项；核心毛利低于10.5\%；重大延期、撤单、客户融资或质保损失；OCF为负且存货/合同资产快于收入；或上市公司公告显示关键订单/高端能力主要位于上市体外。

\begin{riskbox}[硬失效条件]
任一硬条件触发即重开模型：H1归母/扣非低于92/90亿元；核心业务毛利低于10.5\%；年度交付目标实质下修；重大客户履约、质量、制裁或融资事件；全年OCF/归母低于0.8倍且无次期回收路径；主体边界或股本发生重大变化。
\end{riskbox}

\section{最终投委会结论}

中国船舶的基本面优势明确：合并后平台能力完整、在手订单充足、2026利润弹性强。风险也同样明确：订单向毛利和现金的转换缺少细颗粒披露，卖方观点拥挤，当前价已接近独立基础价值。二者合并后的正确动作仍是中性偏多（持有/等待验证），最终目标33.07元、公允区间31.64--35.96元。

\begin{keyinsight}[最后的“所以呢”]
持有是因为订单和盈利方向已经得到A1披露支持；等待是因为+0.15\%目标回报不足以覆盖持续期、现金和执行风险。正式中报若同时提高盈利分母与质量，可快速升级；若只重复订单故事或仅兑现预告，则不改变动作。
\end{keyinsight}

所有后续更新都必须继续遵守四条边界：2025年同一控制下重述口径；2026H1预告不机械年化；中远约500亿元项目已在聚合订单内不重复；军品、沪东中华、集团未分配订单和弱目标价全部零计价。
